\chapter{Trees}
%%%%%%%%%%%%%%%%%%%%%%%%%%%%%%%%%%%%%%%%%%%%%%%%%%%%%%%%%%%%%%%%%%%%%%%%%%%%%%%
\section{Trees}

% Generated by scripts/blueprint_restate.py from TCSlib.GraphTheory.Core.Trees.
% Each body is the STATEMENT only, taken from the declaration's Lean docstring:
% the one-line summary plus the verbatim book statement.  Proof, proof plan,
% significance commentary and formalisation notes are excluded by construction.

\begin{definition}[An edge cut set: a set of edges whose deletion disconnects $G$]
\lean{SimpleGraph.IsEdgeCutSet}
\label{SimpleGraph.IsEdgeCutSet}
\leanok
An \textbf{edge cut set}: a set of edges whose deletion disconnects \texttt{G}.

For subsets \texttt{S} and \texttt{S'} of \texttt{V}, \texttt{[S, S']} denotes the set
of edges with one end in \texttt{S} and the other in \texttt{S'}. An \textbf{edge cut}
of \texttt{G} is a subset of \texttt{E} of the form \texttt{[S, $\bar{S}$]}, where
\texttt{S} is a nonempty proper subset of \texttt{V} and \texttt{$\bar{S}$ = V
\textbackslash{} S}.
\end{definition}

\begin{definition}[A bond (\S{}2.2): a minimal (nonempty) edge cut, i.e]
\lean{SimpleGraph.IsBond}
\label{SimpleGraph.IsBond}
\leanok
\uses{SimpleGraph.IsEdgeCutSet}
A \textbf{bond} (\S{}2.2): a minimal (nonempty) edge cut, i.e. a minimal set of edges
whose
deletion disconnects \texttt{G}.

\emph{A minimal nonempty edge cut of \texttt{G} is called a
bond; each cut edge \texttt{e}, for instance, gives rise to a bond \texttt{\{e\}}. If
\texttt{G} is
connected, then a bond \texttt{B} of \texttt{G} is a minimal subset of \texttt{E} such
that \texttt{G - B} is
disconnected.}
\end{definition}

\begin{definition}[The number of spanning trees $\tau(G)$: the number of tree subgraphs $T \le G$ on the full vertex set]
\lean{SimpleGraph.numSpanningTrees}
\label{SimpleGraph.numSpanningTrees}
\leanok
The \textbf{number of spanning trees} \texttt{$\tau$(G)}: the number of tree subgraphs
\texttt{T $\le$ G} on the
full vertex set.

\emph{We denote the number of spanning trees of \texttt{G} by
\texttt{$\tau$(G)}.}
\end{definition}

\begin{definition}[contract]
\lean{SimpleGraph.contract}
\label{SimpleGraph.contract}
\textbf{Edge contraction} \texttt{G $\cdot$ e} (\S{}2.4). Here it is \emph{stubbed on
the same carrier} \texttt{V} as a placeholder, purely so that the deletion--contraction
recurrence (Theorem 2.8) can be typed.

\emph{An edge \texttt{e} of \texttt{G} is said to be contracted if it is
deleted and its ends are identified; the resulting graph is denoted \texttt{G $\cdot$
e}.}
\end{definition}

\begin{theorem}[Unique path between vertices in a tree]
\lean{SimpleGraph.tree_unique_path}
\label{SimpleGraph.tree_unique_path}
\leanok
Let $G$ be a tree on a finite vertex set, and let $u$ and $v$ be any two of its
vertices. Then there is exactly one path in $G$ from $u$ to $v$; that is, among all
walks from $u$ to $v$, precisely one has no repeated vertices.
\end{theorem}

\begin{theorem}[Edge count of a finite tree]
\lean{SimpleGraph.tree_card_edgeFinset}
\label{SimpleGraph.tree_card_edgeFinset}
\leanok
Let $G$ be a finite simple graph that is a tree, meaning it is connected and contains no
cycles. Then the number of edges of $G$ is one less than the number of vertices; that
is, if $G$ has $n$ vertices, then it has exactly $n - 1$ edges.
\end{theorem}

\begin{theorem}[Every nontrivial finite tree has at least two leaves]
\lean{SimpleGraph.tree_two_leaves}
\label{SimpleGraph.tree_two_leaves}
\leanok
Let $G$ be a tree on a finite vertex set $V$ with at least two vertices. Then $G$ has at
least two leaves: there exist distinct vertices $u, v \in V$ with $\deg(u) = 1$ and
$\deg(v) = 1$.
\end{theorem}

\begin{theorem}[Characterization of bridges by cycles]
\lean{SimpleGraph.cutEdge_iff_no_cycle}
\label{SimpleGraph.cutEdge_iff_no_cycle}
\leanok
Let $G$ be a finite simple graph and let $v,w$ be two of its vertices. The edge
$\{v,w\}$ is a bridge of $G$ — that is, its removal increases the number of connected
components — if and only if $v$ and $w$ are adjacent in $G$ and the edge $\{v,w\}$
occurs in no cycle of $G$.
\end{theorem}

\begin{theorem}[Trees as graphs whose edges are all bridges]
\lean{SimpleGraph.connected_isTree_iff_forall_edge_isBridge}
\label{SimpleGraph.connected_isTree_iff_forall_edge_isBridge}
\leanok
Let $G$ be a connected finite simple graph. Then $G$ is a tree — that is, connected and
acyclic — if and only if every edge of $G$ is a bridge, meaning that deleting it
disconnects the graph (leaving its two endpoints in different connected components).
\end{theorem}

\begin{theorem}[Existence of a spanning tree]
\lean{SimpleGraph.exists_spanningTree}
\label{SimpleGraph.exists_spanningTree}
\leanok
Let $G$ be a connected simple graph on a finite vertex set $V$. Then there is a spanning
subgraph $T$ of $G$ — a simple graph on the same vertex set $V$ with $T \le G$ — that is
a tree, that is, $T$ is connected and contains no cycle.
\end{theorem}

\begin{theorem}[Connected graphs have at least $n-1$ edges]
\lean{SimpleGraph.connected_card_edgeFinset_ge}
\label{SimpleGraph.connected_card_edgeFinset_ge}
\leanok
Let $G$ be a connected simple graph on a finite vertex set $V$ with $|V| = n$. Then the
number of edges of $G$ is at least $n - 1$.
\end{theorem}

\begin{theorem}[Adding a non-tree edge creates a cycle]
\lean{SimpleGraph.spanningTree_add_edge_exists_cycle}
\label{SimpleGraph.spanningTree_add_edge_exists_cycle}
\leanok
Let $G$ be a finite simple graph, and let $T$ be a spanning tree of $G$: a connected,
acyclic subgraph of $G$ that touches every vertex. If $e$ is an edge of $G$ that does
not lie in $T$, then the graph $T + e$ obtained by adjoining $e$ to $T$ contains a
cycle, i.e. there is a vertex $u$ and a closed walk from $u$ to itself in $T + e$ that
is a cycle.
\end{theorem}

\begin{theorem}[The unique bond in a cotree with one tree edge]
\lean{SimpleGraph.cotree_bond}
\label{SimpleGraph.cotree_bond}
\leanok
\uses{SimpleGraph.IsBond}
Let $G$ be a finite connected graph, let $T$ be a spanning tree of $G$, and let $e$ be
an edge of $T$. Write $\bar T = E(G) \setminus E(T)$ for the cotree, the set of edges of
$G$ that lie outside $T$. Then (i) no bond of $G$ is contained in $\bar T$, and (ii)
there is exactly one bond of $G$ contained in $\{e\} \cup \bar T$.
\end{theorem}

\begin{theorem}[Cut vertices of a tree are the non-leaf vertices]
\lean{SimpleGraph.tree_isCutVertex_iff_degree}
\label{SimpleGraph.tree_isCutVertex_iff_degree}
\leanok
\uses{SimpleGraph.IsCutVertex}
Let $G$ be a tree, that is, a finite connected acyclic graph, and let $v$ be a vertex of
$G$. Then $v$ is a cut vertex of $G$ — meaning $G$ is connected and deleting $v$
(together with its incident edges) leaves a disconnected graph — if and only if the
degree of $v$ satisfies $\deg(v) > 1$.
\end{theorem}

\begin{theorem}[Every connected graph has two non-cut vertices]
\lean{SimpleGraph.connected_two_non_cutVertices}
\label{SimpleGraph.connected_two_non_cutVertices}
\leanok
\uses{SimpleGraph.IsCutVertex}
Let $G$ be a finite connected simple graph with at least two vertices. Then $G$ has two
distinct vertices $u$ and $v$, neither of which is a cut vertex.
\end{theorem}

\begin{theorem}[Deletion–contraction for spanning trees]
\lean{SimpleGraph.numSpanningTrees_deletion_contraction}
\label{SimpleGraph.numSpanningTrees_deletion_contraction}
\uses{SimpleGraph.contract, SimpleGraph.numSpanningTrees}
Let $G$ be a finite graph, and write $\tau(G)$ for its number of spanning trees, that
is, the number of tree subgraphs of $G$ on the full vertex set. If $e$ is an edge of
$G$, then
\[
\tau(G) = \tau(G - e) + \tau(G \cdot e),
\]
where $G - e$ is the graph obtained by deleting $e$ and $G \cdot e$ is the graph
obtained by contracting $e$ (deleting it and identifying its two ends).
\end{theorem}

\begin{theorem}[Cayley's formula for the number of spanning trees of the complete graph]
\lean{SimpleGraph.cayley}
\label{SimpleGraph.cayley}
\leanok
\uses{SimpleGraph.numSpanningTrees}
For every natural number $n$, the number of spanning trees of the complete graph $K_n$
on $n$ vertices is $\tau(K_n) = n^{n-2}$, where the exponent is the truncated difference
(so it is $0$ when $n < 2$).
\end{theorem}

\begin{theorem}[Existence of a minimum-weight spanning tree]
\lean{SimpleGraph.exists_min_weight_spanningTree}
\label{SimpleGraph.exists_min_weight_spanningTree}
\leanok
Let $G$ be a connected simple graph on a finite vertex set $V$, and let $w$ assign a
real weight $w(e)$ to each unordered pair $e$ of vertices. Then there exists a spanning
subgraph $T$ of $G$ that is a tree and whose total edge weight $\sum_{e \in E(T)} w(e)$
is least among all subgraphs $T'$ of $G$ that are trees.
\end{theorem}
